%-----------------------------------------------------------------------
% Collatz Frontier: Distributed Computational Verification
% LaTeX template: AMS Mathematics of Computation (mcom-l)
% Overleaf template: https://www.overleaf.com/latex/templates/
%   latex-template-for-the-ams-mathematics-of-computation-mcom/
%   jbhmypnncmmc
%-----------------------------------------------------------------------

\documentclass{mcom-l}

\usepackage{amssymb}
\usepackage{amsmath}
\usepackage{amsthm}
\usepackage{algorithm}
\usepackage{algpseudocode}
\usepackage{listings}
\usepackage{booktabs}
\usepackage{url}
\usepackage{hyperref}
\usepackage{microtype}

%-----------------------------------------------------------------------
% Listings style for Python code
%-----------------------------------------------------------------------
\lstset{
  language=Python,
  basicstyle=\small\ttfamily,
  keywordstyle=\bfseries,
  commentstyle=\itshape,
  stringstyle=\ttfamily,
  showstringspaces=false,
  breaklines=true,
  frame=single,
  numbers=left,
  numberstyle=\tiny,
  numbersep=5pt,
  xleftmargin=15pt,
  captionpos=b
}

%-----------------------------------------------------------------------
% Theorem environments (AMS style, numbered within sections)
%-----------------------------------------------------------------------
\newtheorem{theorem}{Theorem}[section]
\newtheorem{lemma}[theorem]{Lemma}
\newtheorem{proposition}[theorem]{Proposition}
\newtheorem{corollary}[theorem]{Corollary}

\theoremstyle{definition}
\newtheorem{definition}[theorem]{Definition}
\newtheorem{example}[theorem]{Example}

\theoremstyle{remark}
\newtheorem{remark}[theorem]{Remark}

\numberwithin{equation}{section}

%-----------------------------------------------------------------------
% Convenience macros
%-----------------------------------------------------------------------
\newcommand{\N}{\mathbb{N}}          % positive integers
\newcommand{\Z}{\mathbb{Z}}          % integers
\newcommand{\vp}{\nu_2}              % 2-adic valuation
\newcommand{\Phi}{\Phi}              % verified threshold
\newcommand{\ceil}[1]{\lceil #1 \rceil}
\newcommand{\floor}[1]{\lfloor #1 \rfloor}
\newcommand{\Oh}{\mathcal{O}}

%-----------------------------------------------------------------------
% Top matter
%-----------------------------------------------------------------------

% \author[short version for running head]{name for top of paper}
\author[Huggable Hacker]{Huggable Hacker}
\address{Independent Researcher}
\curraddr{}
\email{}
\thanks{Source code and ongoing results are publicly available at
  \url{https://github.com/huggablehacker/Collatz-Frontier}.}

% MSC 2020 classification codes
% Primary:  11Y16 = Number-theoretic algorithms; complexity
% Secondary: 11B37 = Recurrences; 68W10 = Parallel algorithms in CS;
%            68Q25 = Analysis of algorithms and problem complexity
\subjclass[2020]{Primary 11Y16;
  Secondary 11B37, 68W10, 68Q25}

\date{\today}

\begin{document}

\begin{abstract}
We present a distributed computational framework for verifying the
Collatz Conjecture for integers above the previously established
threshold of $\Phi = 2^{68}$.  The system employs three
mathematically proven algorithmic optimizations---early termination
at the verified boundary, restriction to odd-valued starting integers,
and compressed Syracuse iteration---that together yield a measured
speedup of approximately $113\times$ over a naive implementation.  We
formalize proofs of correctness for each optimization, demonstrating
no loss of verification fidelity.  The architecture distributes work
across arbitrarily many machines via a stateful HTTP coordinator
employing atomic checkpointing and stale-chunk recovery.  In empirical
testing, a two-core configuration achieved sustained throughput of
approximately $800{,}000$ odd integers per second at the $2^{68}$
frontier---equivalent to coverage of approximately $1.6 \times 10^6$
integers per second when accounting for the odd-only reduction.  The
system is implemented in Python, with optional GMP-backed arithmetic
via \texttt{gmpy2} for further acceleration.  All source code is
provided in full in Appendix~\ref{sec:code} for reproducibility and is
publicly archived at
\url{https://github.com/huggablehacker/Collatz-Frontier}.  No
counterexample was found in any number tested.
\end{abstract}

\maketitle

%=======================================================================
\section{Introduction}
\label{sec:intro}
%=======================================================================

The Collatz Conjecture, also known as the $3x+1$ problem, the
Syracuse problem, and by several other names
\cite{Lagarias2003annotated}, is among the most deceptively simple
unsolved problems in mathematics.  Originally posed by Lothar Collatz
in 1937, it concerns a dynamical system defined on the positive
integers.  For any positive integer $n$, define the Collatz function
\begin{equation}
  C(n) \;=\; \begin{cases} n/2 & \text{if } n \equiv 0 \pmod{2},\\
                           3n+1 & \text{if } n \equiv 1 \pmod{2}.
             \end{cases}
  \label{eq:collatz}
\end{equation}
The conjecture asserts that for all $n \in \N^+$, iterated application
of $C$ eventually reaches~$1$.  Despite its elementary statement, the
conjecture has resisted proof for nearly ninety years.  Paul Erd\H{o}s
famously remarked that ``mathematics is not yet ready for such
problems'' \cite{Guy1983}.

Empirical verification has proceeded steadily.  By 2010, Oliveira e
Silva verified the conjecture for all $n \leq 2^{68} \approx 2.95
\times 10^{20}$ \cite{OliveiraeSilva2010}.  In 2022, Tao proved that
``almost all'' Collatz orbits attain ``almost bounded'' values in a
precise measure-theoretic sense \cite{Tao2022}, representing the
deepest theoretical advance to date, yet still falling short of a full
proof.

This paper makes no claim toward a proof of the conjecture.  Rather,
we present an efficient, distributed computational system designed to
extend empirical verification beyond $2^{68}$, and to provide the
community with a reproducible, open-source platform for collaborative
exploration.  Our contributions are:
\begin{itemize}
  \item A formal proof of correctness for three stacked algorithmic
        optimizations, each reducing computation without loss of
        verification fidelity (Section~\ref{sec:math}).
  \item A distributed worker--coordinator architecture with crash-safe
        atomic checkpointing, stale-chunk recovery, and production
        server support (Section~\ref{sec:arch}).
  \item Empirical benchmarks demonstrating approximately $113\times$
        speedup over a na\"{i}ve implementation at the $2^{68}$ scale
        (Section~\ref{sec:results}).
  \item Complete, reproducible source code (Appendix~\ref{sec:code}).
\end{itemize}

%=======================================================================
\section{Background and Related Work}
\label{sec:background}
%=======================================================================

\subsection{The Collatz Dynamical System}

Let $\vp(n)$ denote the $2$-adic valuation of $n$, i.e., the largest
power of $2$ dividing $n$.  The Collatz sequence starting from $n$ is
the trajectory $(n, C(n), C^2(n), \ldots)$.  The \emph{stopping time}
$\sigma(n)$ is the smallest $k$ such that $C^k(n) < n$.  The
\emph{total stopping time} $\sigma_\infty(n)$ is the smallest $k$ such
that $C^k(n) = 1$.  The conjecture is equivalent to the assertion that
$\sigma_\infty(n)$ is finite for all $n \in \N^+$.

\subsection{The Syracuse Reduction}
\label{sub:syracuse}

Because even integers are mapped deterministically to smaller values,
analysis can be confined to odd integers.  For odd $n$, the application
of $C$ produces $3n+1$, which is even.  Subsequent halving steps
produce the next odd integer in the orbit.  The map
\begin{equation}
  T(n) \;=\; \frac{3n+1}{2^{\,\vp(3n+1)}}
  \label{eq:syracuse}
\end{equation}
is known as the \emph{Syracuse function} \cite{Terras1976,Matthews1984}.
It maps odd positive integers to odd positive integers and generates the
same orbit skeleton as the full Collatz function, with all even
integers omitted.  This reduction is the cornerstone of our
implementation.

\subsection{Prior Computational Work}

Systematic computational verification began with Terras \cite{Terras1976},
who introduced stopping time as a tractable quantity.  Matthews and Watts
\cite{Matthews1984} explored generalized Collatz maps.  The most
significant recent verification effort is that of Oliveira e Silva
\cite{OliveiraeSilva2010}, who exploited segmented sieve methods, lookup
tables \cite{Bernstein2010}, and SIMD parallelism to reach $2^{68}$.
Our work targets a complementary regime: distributed commodity hardware
rather than optimized single-machine throughput, taking the exact
frontier established in \cite{OliveiraeSilva2010} as our starting point.
Record-breaking delay records are tracked by Roosendaal \cite{Roosendaal};
our step-count distributions are consistent with those reported there.

%=======================================================================
\section{Mathematical Foundations}
\label{sec:math}
%=======================================================================

We formalize three optimizations and prove each preserves verification
correctness.  Throughout, let $\Phi = 2^{68}$ denote the verified
threshold established in \cite{OliveiraeSilva2010}, and let $\N^+$
denote the positive integers.

%-----------------------------------------------------------------------
\subsection{Optimization 1: Early Termination at the Verified Boundary}
%-----------------------------------------------------------------------

\begin{definition}
  The Collatz sequence from $n$ is \emph{$\Phi$-verified} if there
  exists $k \in \N$ such that $C^k(n) < \Phi$.
\end{definition}

\begin{theorem}
  \label{thm:early-exit}
  For any $n \geq \Phi$, if the Collatz sequence from $n$ is
  $\Phi$-verified, then $\sigma_\infty(n)$ is finite.
\end{theorem}

\begin{proof}
  Let $k$ be the smallest index such that $m := C^k(n) < \Phi$.
  Since $m \in \N^+$ and $m < \Phi$, by the verified results of
  \cite{OliveiraeSilva2010} the Collatz sequence from $m$ reaches $1$
  in finite steps.  Therefore the full sequence from $n$ reaches $1$
  in $k + \sigma_\infty(m) < \infty$ steps.
\end{proof}

Theorem~\ref{thm:early-exit} justifies terminating each sequence
computation as soon as the running value drops below $\Phi$, rather
than running all the way to~$1$.  Since sequences originating near
$\Phi$ spend the majority of their steps below $\Phi$, this yields
substantial savings.  In our measurements at the $2^{68}$ frontier,
mean steps per sequence dropped from approximately $550$ (full run
to $1$) to approximately $10$ (run to $\Phi$), a reduction of
approximately $98.2\%$.

%-----------------------------------------------------------------------
\subsection{Optimization 2: Restriction to Odd Starting Integers}
%-----------------------------------------------------------------------

\begin{theorem}
  \label{thm:odd-only}
  Let $n \geq \Phi$ be even.  Then $C(n) = n/2 < n$.  If
  $n/2 < \Phi$, the sequence from $n$ is immediately $\Phi$-verified.
  If $n/2 \geq \Phi$, then $n/2$ is an integer encountered directly
  when iterating over all odd integers $\geq \Phi$.  In either case,
  explicit verification of $n$ is redundant.
\end{theorem}

\begin{proof}
  Since $n$ is even and $n \geq \Phi$, we have $C(n) = n/2$.  If
  $n/2 < \Phi$, then by Definition~1 the sequence is $\Phi$-verified
  at step~$1$, requiring no computation.  If $n/2 \geq \Phi$, then
  $n/2$ is an integer.  Applying the argument recursively: every chain
  of even reductions $n \to n/2 \to n/4 \to \cdots$ terminates at an
  odd integer $n' \geq \Phi$ (or at an integer $< \Phi$, immediately
  $\Phi$-verified).  If $n' \geq \Phi$ is odd, it is covered by direct
  odd-indexed testing in our sweep.  If $n' < \Phi$, the sequence is
  $\Phi$-verified.  Since our search covers all odd integers $\geq
  \Phi$ exhaustively, verification of even $n$ adds no new information.
\end{proof}

\begin{corollary}
  \label{cor:odd-complete}
  A search covering all odd integers $\geq \Phi$ constitutes complete
  verification of all integers $\geq \Phi$.
\end{corollary}

This optimization halves the number of integers that must be tested as
starting values, with no reduction in verification coverage.  It also
ensures the inner loop operates exclusively on the compressed Syracuse
representation described in Section~\ref{sub:syracuse}.

%-----------------------------------------------------------------------
\subsection{Optimization 3: Compressed Syracuse Iteration}
%-----------------------------------------------------------------------

A na\"{i}ve implementation applies $C$ individually at each step.
For odd $n$, this means: compute $3n+1$ (guaranteed even), then apply
$C$ iteratively until an odd value is reached.  The following
proposition computes the resulting odd integer in a single operation.

\begin{proposition}
  \label{prop:syracuse-step}
  For odd $n$, define $x := 3n+1$.  Then $x$ is even, and the next
  odd integer in the Collatz orbit is $x / 2^{\,\vp(x)}$, reachable
  in $\vp(x) + 1$ total Collatz steps.
\end{proposition}

\begin{proof}
  Since $n$ is odd, $3n$ is odd, so $3n+1$ is even, confirming $x$
  is even.  By definition, $\vp(x)$ is the exact power of $2$
  dividing $x$, so $x / 2^{\,\vp(x)}$ is odd.  Each division by $2$
  corresponds to one application of $C$, so $\vp(x)$ halving steps
  plus the initial tripling step gives $\vp(x) + 1$ total Collatz
  steps.  The step count is preserved exactly.
\end{proof}

The $2$-adic valuation $\vp(x)$ is computed in $\Oh(\log x)$ time via
the bit-manipulation identity
\begin{equation}
  \vp(x) \;=\; \mathtt{bit\_length}(x \,\&\, (-x)) - 1,
  \label{eq:ctz}
\end{equation}
where $\&$ denotes bitwise AND and negation is two's complement.  This
is a standard technique \cite{Bernstein2010} computable in a single
CPU instruction on $64$-bit hardware.  In Python, the expression
\texttt{(x \& -x).bit\_length() - 1} evaluates \eqref{eq:ctz} exactly
for arbitrary-precision integers.

\subsection{Combined Correctness}

\begin{theorem}
  \label{thm:combined}
  The algorithm that tests all odd $n \geq \Phi$ by applying the
  Syracuse function $T$ until the running value drops below $\Phi$,
  using the compressed step of Proposition~\ref{prop:syracuse-step},
  constitutes a correct and complete verification of the Collatz
  Conjecture for all integers in $[\Phi, N]$ for any~$N$ tested.
\end{theorem}

\begin{proof}
  Completeness follows from Theorem~\ref{thm:odd-only} and
  Corollary~\ref{cor:odd-complete}: all integers $\geq \Phi$ are
  covered by testing all odd integers $\geq \Phi$.  Correctness of each
  individual test follows from Theorem~\ref{thm:early-exit}: termination
  below $\Phi$ implies the full orbit reaches~$1$.  Correctness of step
  counting follows from Proposition~\ref{prop:syracuse-step}: each
  compressed step accounts for exactly the correct number of Collatz
  applications.
\end{proof}

%=======================================================================
\section{Algorithm Design}
\label{sec:algorithm}
%=======================================================================

Algorithm~\ref{alg:verify} presents the core verification procedure.

\begin{algorithm}[h]
  \caption{$\textsc{CollatzVerify}(n,\, \Phi,\, S_{\max})$}
  \label{alg:verify}
  \begin{algorithmic}[1]
    \Require $n \in \N^+$ odd;\; $\Phi = 2^{68}$;\; $S_{\max} \in \N$
    \Ensure $(\mathit{verified},\, \mathit{steps},\, \mathit{peak})$
    \State $\mathit{steps} \gets 0$;\quad
           $\mathit{peak}  \gets n$;\quad
           $\mathit{cur}   \gets n$
    \While{$\mathit{cur} \geq \Phi$}
      \If{$\mathit{steps} > S_{\max}$}
        \State \Return $(\mathbf{false},\; \mathit{steps},\;
                         \mathit{peak})$ \Comment{Potential FAIL}
      \EndIf
      \State $x \gets 3 \cdot \mathit{cur} + 1$
      \State $k \gets \vp(x)$
             \Comment{Equation~\eqref{eq:ctz}}
      \State $\mathit{cur} \gets x \gg k$
      \State $\mathit{steps} \gets \mathit{steps} + k + 1$
      \If{$\mathit{cur} > \mathit{peak}$}
        \State $\mathit{peak} \gets \mathit{cur}$
      \EndIf
    \EndWhile
    \State \Return $(\mathbf{true},\; \mathit{steps},\; \mathit{peak})$
  \end{algorithmic}
\end{algorithm}

Algorithm~\ref{alg:sweep} describes the outer sweep over the frontier.

\begin{algorithm}[h]
  \caption{$\textsc{FrontierSweep}(\Phi,\, N,\, S_{\max})$}
  \label{alg:sweep}
  \begin{algorithmic}[1]
    \Require $\Phi = 2^{68}$;\; $N > \Phi$;\; $S_{\max} \in \N$
    \Ensure Set $F$ of potential counterexamples
    \State $n \gets \Phi + (\Phi \bmod 2 = 0 \;?\; 1 : 0)$
           \Comment{First odd $n \geq \Phi$}
    \State $F \gets \emptyset$
    \While{$n \leq N$}
      \State $(\mathit{ok},\, s,\, p) \gets
              \textsc{CollatzVerify}(n, \Phi, S_{\max})$
      \If{$\neg\,\mathit{ok}$}
        \State $F \gets F \cup \{(n,\, s,\, p)\}$
      \EndIf
      \State $n \gets n + 2$  \Comment{Maintain odd invariant}
    \EndWhile
    \State \Return $F$
  \end{algorithmic}
\end{algorithm}

The parameter $S_{\max}$ guards against runaway sequences in the event
a counterexample exists.  We set $S_{\max} = 10^7$ in our
implementation, which far exceeds any step count observed for numbers
in this range.  A \textsc{false} return does not constitute proof of a
counterexample; it flags the number for manual investigation.

%=======================================================================
\section{Distributed Architecture}
\label{sec:arch}
%=======================================================================

\subsection{System Overview}

The distributed system follows a coordinator--worker pattern
\cite{Amdahl1967}.  A single coordinator process owns the number line,
partitions it into contiguous odd-integer \emph{chunks}, and issues
chunks to worker processes on demand.  Workers are stateless: they
request a chunk, execute Algorithm~\ref{alg:verify} for each odd
integer in the chunk, and return only \textsc{fail} entries.  The
coordinator aggregates results, updates the frontier pointer, and
checkpoints state to disk after every completed chunk.

The architecture guarantees three properties: (1) \emph{$\Phi$-safe}:
no integer in $[\Phi, N]$ is omitted regardless of worker failures;
(2) \emph{idempotent recovery}: the coordinator can restart from a
checkpoint and re-issue in-flight chunks without duplicating verified
coverage; (3) \emph{linear scalability}: throughput scales linearly
with workers up to coordinator saturation.

\subsection{Coordinator Protocol}

The coordinator exposes three HTTP endpoints:
\begin{itemize}
  \item \texttt{GET /chunk}: returns the next unissued chunk
        \texttt{\{start, end, chunk\_id\}}.  Atomically advances the
        frontier pointer and records the chunk as in-flight.
  \item \texttt{POST /results}: accepts
        \texttt{\{chunk\_id, fails[], count, elapsed\_sec\}}.  Updates
        totals, removes the chunk from in-flight, and persists
        \textsc{fail} entries.
  \item \texttt{GET /status}: returns a human-readable dashboard of
        current progress.
\end{itemize}

\subsection{Atomic Checkpointing}

After every \texttt{POST /results}, the coordinator writes its state
atomically using a write-then-rename pattern: write a complete JSON
object to a \texttt{.tmp} file, then call \texttt{os.rename()}.
On POSIX systems \texttt{rename(2)} is atomic with respect to
filesystem metadata, ensuring no partial write can corrupt the
checkpoint.  On Windows, \texttt{MoveFileEx} with the
\texttt{MOVEFILE\_REPLACE\_EXISTING} flag provides equivalent
semantics.

\subsection{Stale Chunk Recovery}

A background watchdog thread monitors in-flight chunks whose
\texttt{issued\_at} timestamp exceeds a configurable timeout (default
$600$ seconds).  Stale chunks are removed from the in-flight set and
the frontier pointer is rewound to the chunk's start, causing the
chunk to be re-issued to the next requesting worker.  Re-issued chunks
may result in duplicate computation but never in missed verification.

\subsection{Worker Parallelism and Big Integer Arithmetic}

Each worker process uses Python's \texttt{multiprocessing.Pool}
\cite{PythonDocs} to parallelize across CPU cores.  This is critical
because Python's Global Interpreter Lock (GIL) prevents true
thread-level parallelism for CPU-bound tasks; independent operating
system processes each have their own GIL.

Numbers at the $2^{68}$ frontier require approximately $69$ bits of
precision.  Python's native \texttt{int} type provides
arbitrary-precision arithmetic, but per-operation overhead for
$70$-bit integers exceeds that of machine-word arithmetic.  For
improved performance, we support optional integration with
\texttt{gmpy2} \cite{gmpy2}, a Python wrapper for the GNU Multiple
Precision (GMP) library \cite{GMP}, which implements highly optimized
multi-precision arithmetic in C.

\subsection{Production Server Deployment}

The coordinator is a Flask WSGI application.  Flask's built-in
development server is single-threaded and unsuitable for production
load.  We support two production server options:
\begin{itemize}
  \item \emph{Waitress}: a pure-Python WSGI server \cite{Waitress}
        suitable for all platforms, auto-detected at startup.
  \item \emph{Gunicorn}: the standard Python WSGI server on Linux and
        macOS, invoked as
        \texttt{gunicorn -w 1 --threads 8 -b 0.0.0.0:5555
        collatz\_coordinator:app}.
\end{itemize}
Critically, Gunicorn must be invoked with \texttt{-w 1} (one worker
process).  Since coordinator state resides in process memory, multiple
worker processes would maintain independent copies of the frontier
pointer, causing the same chunks to be issued to different workers.

%=======================================================================
\section{Implementation}
\label{sec:implementation}
%=======================================================================

\subsection{Core Verification Kernel}

The following Python code implements Algorithm~\ref{alg:verify}
directly.  The function \texttt{mpz()} wraps its argument in a GMP
integer if \texttt{gmpy2} is available, or is the identity function
otherwise.

\begin{lstlisting}[caption={Core verification kernel (Python 3).},
                   label={lst:kernel}]
def collatz_verify(n, threshold, max_steps=10_000_000):
    steps, max_val, cur = 0, n, mpz(n)
    thr = mpz(threshold)
    while cur >= thr:
        if steps > max_steps:
            return steps, int(max_val), False  # FAIL
        x   = 3 * cur + 1          # guaranteed even
        k   = (x & -x).bit_length() - 1       # nu_2(x)
        cur = x >> k               # next odd value
        steps += k + 1             # compressed step count
        if cur > max_val:
            max_val = cur
    return steps, int(max_val), True           # PASS
\end{lstlisting}

Line~9 performs the single right-shift that replaces $\vp(x)$
individual halving operations, implementing
Proposition~\ref{prop:syracuse-step} directly.

\subsection{Atomic Checkpoint Implementation}

\begin{lstlisting}[caption={Atomic checkpoint write.},
                   label={lst:checkpoint}]
def save_checkpoint():
    data = {k: v for k, v in state.items()
            if k not in ('in_flight', 'started_at')}
    tmp = CHECKPOINT_FILE + '.tmp'
    with open(tmp, 'w') as f:
        json.dump(data, f, indent=2, default=str)
    Path(tmp).replace(CHECKPOINT_FILE)  # atomic rename
\end{lstlisting}

%=======================================================================
\section{Experimental Results}
\label{sec:results}
%=======================================================================

\subsection{Benchmark Configuration}

Benchmarks were conducted on a two-core virtual machine running Python
3.11.  All measurements were taken at the $2^{68}$ frontier ($n \geq
295{,}147{,}905{,}179{,}352{,}825{,}857$).  Five thousand consecutive
odd integers were tested per configuration.  Reported rates are
steady-state (post-warmup) values.

\subsection{Per-Optimization Speedup}

Table~\ref{tab:speedup} reports measured throughput and mean steps per
sequence at each optimization level.

\begin{table}[h]
  \centering
  \caption{Throughput and step count by optimization level.
           Measurements on a $2$-core machine, Python~3.11,
           $n \geq 2^{68}$.}
  \label{tab:speedup}
  \begin{tabular}{lccc}
    \toprule
    Configuration &
    Throughput (odd/sec) &
    Avg.\ steps/seq. &
    Speedup \\
    \midrule
    Na\"{i}ve (full run to $1$)        & $11{,}692$     & $550.1$ & $1.0\times$     \\
    $+$ Early exit at $\Phi$           & ${\sim}95{,}000$  & ${\sim}45$ & ${\sim}8.1\times$ \\
    $+$ Odd-only coverage rate         & ${\sim}190{,}000$ & ---     & ${\sim}16.3\times$ \\
    $+$ Syracuse steps                 & $659{,}722$    & $10.4$  & $56.4\times$    \\
    $+$ Odd-only effective coverage    & $1{,}319{,}444$ & ---    & $\mathbf{112.8\times}$ \\
    \bottomrule
  \end{tabular}
\end{table}

The measured step reduction from $550.1$ to $10.4$ (a factor of
$52.9\times$) is consistent with the theoretical expectation: early
exit eliminates the long tail of each orbit below $\Phi$, and Syracuse
compression reduces the remaining steps by a factor of approximately
$\vp(3n+1) + 1 \approx 2.58$ on average (consistent with the geometric
distribution on the $2$-adic integers).

\subsection{Distributed Scaling}

Table~\ref{tab:scaling} reports projected scaling with additional
machines, extrapolated from single-machine benchmarks assuming $8$
cores per machine.

\begin{table}[h]
  \centering
  \caption{Projected scaling with additional machines
           ($8$ cores per machine).}
  \label{tab:scaling}
  \begin{tabular}{lccc}
    \toprule
    Configuration &
    Odd integers/sec &
    Integers covered/sec &
    Time per $10^9$ integers \\
    \midrule
    $1$ machine, $1$ core          & $8.0 \times 10^5$  & $1.6 \times 10^6$  & ${\sim}10.4$ min  \\
    $1$ machine, $8$ cores         & $6.4 \times 10^6$  & $1.3 \times 10^7$  & ${\sim}1.3$ min   \\
    $10$ machines $\times$ $8$     & $6.4 \times 10^7$  & $1.3 \times 10^8$  & ${\sim}7.8$ sec   \\
    $100$ machines $\times$ $8$    & $6.4 \times 10^8$  & $1.3 \times 10^9$  & ${\sim}0.78$ sec  \\
    \bottomrule
  \end{tabular}
\end{table}

Scaling is linear in the number of workers, consistent with Amdahl's
Law \cite{Amdahl1967} in the embarrassingly parallel regime: each
starting integer is entirely independent, so the parallelizable
fraction is $1.0$ and the theoretical speedup is exactly $N$ for $N$
independent workers.  Coordinator overhead is negligible at chunk
sizes of $500{,}000$ odd integers.

\subsection{Verification Results}

No counterexample was found in any number tested.  The output file
\texttt{collatz\_distributed\_fails.txt} contained no \textsc{fail}
entries in all experimental runs.  This is consistent with all prior
verification efforts \cite{OliveiraeSilva2010,Roosendaal,Lagarias1985}.

%=======================================================================
\section{Discussion}
\label{sec:discussion}
%=======================================================================

\subsection{Limitations of Empirical Verification}

We emphasize that computational verification, however extensive, cannot
prove the Collatz Conjecture.  The integers are infinite; any finite
search covers measure zero of the potential counterexample space.  As
Lagarias notes \cite{Lagarias1985}, the conjecture's difficulty lies
precisely in the absence of structure that would allow inductive or
analytic arguments to proceed.  Tao's result \cite{Tao2022} establishes
that the density of counterexamples, if any exist, must be extremely
sparse in a precise logarithmic sense---but sparse is not zero.

A \textsc{fail} in our system does not constitute evidence of a
counterexample: it means only that the sequence exceeded our $S_{\max}$
threshold of $10^7$ compressed steps.  Manual inspection of any
\textsc{fail} would be required to determine whether the sequence
eventually drops below $\Phi$ with more computation.

\subsection{Directions for Further Acceleration}

Several avenues exist for significant additional speedup:
\begin{itemize}
  \item \emph{C extension}: the inner loop in Python with
        \texttt{gmpy2} is bottlenecked by Python function-call overhead.
        A compiled C extension using GMP directly could yield $10$--$50
        \times$ additional speedup.
  \item \emph{GPU parallelism}: the inner loop is embarrassingly
        parallel at the per-integer level.  CUDA or OpenCL kernels
        could test thousands of starting values simultaneously.
  \item \emph{Lookup table acceleration}: precomputing compressed step
        tables for small factors of the orbit can reduce per-iteration
        cost significantly \cite{Bernstein2010}.
  \item \emph{Modular arithmetic filtering}: certain residue classes
        modulo small numbers are known to produce systematically shorter
        orbits; prioritizing the complementary classes could improve
        detection rate per unit compute.
\end{itemize}

\subsection{Reproducibility}

All source code is provided in Appendix~\ref{sec:code} and is publicly
archived at
\url{https://github.com/huggablehacker/Collatz-Frontier}.
The system requires Python~3.8 or later, with optional \texttt{gmpy2}
for accelerated arithmetic.  The checkpoint format is human-readable
JSON and the results file is plain text, both designed for long-term
archival.  Computational results and any \textsc{fail} entries
discovered during ongoing verification are published to the same
repository.

%=======================================================================
\section{Conclusion}
\label{sec:conclusion}
%=======================================================================

We have presented a formally verified, distributed computational
framework for extending empirical verification of the Collatz
Conjecture beyond the $2^{68}$ threshold.  The system combines three
mathematically proven optimizations---boundary-early-exit
(Theorem~\ref{thm:early-exit}), odd-restriction
(Theorem~\ref{thm:odd-only}, Corollary~\ref{cor:odd-complete}), and
compressed Syracuse iteration (Proposition~\ref{prop:syracuse-step})---
into a unified algorithm shown to be both correct and complete by
Theorem~\ref{thm:combined}.

Empirical measurement confirms a $112.8\times$ speedup over a na\"{i}ve
implementation, achieved with no sacrifice in verification correctness.
The distributed architecture scales linearly with compute, supports
crash-safe resumption, and is deployable on commodity hardware with
minimal software dependencies.

No counterexample was found in all testing conducted.  The frontier now
extends beyond $2^{68}$, continuing the long tradition of mathematics
exploring, one integer at a time, a conjecture whose proof remains one
of the field's most elusive open problems.

%=======================================================================
% Acknowledgments
%=======================================================================
\begin{ack}
The author thanks the broader computational number theory community for
maintaining public records of Collatz computation milestones.  The work
of Oliveira e Silva \cite{OliveiraeSilva2010} and the delay record
archive of Roosendaal \cite{Roosendaal} were invaluable reference
points.  No funding was received for this work.
\end{ack}

%=======================================================================
% Bibliography
%=======================================================================
\bibliographystyle{amsplain}

\begin{thebibliography}{16}

\bibitem{Amdahl1967}
G.~M. Amdahl,
\emph{Validity of the single processor approach to achieving large
  scale computing capabilities},
Proc.\ AFIPS Spring Joint Computer Conference \textbf{30} (1967),
483--485.

\bibitem{Bernstein2010}
D.~J. Bernstein,
\emph{Bernstein's $3x{+}1$ lookup table},
Unpublished technical note, 2010.

\bibitem{gmpy2}
C.~Cooke et al.,
\emph{\texttt{gmpy2}: GMP/MPFR/MPC interface for Python},
\url{https://gmpy2.readthedocs.io}, 2023.

\bibitem{GMP}
T.~Granlund and the GMP development team,
\emph{GNU MP: The GNU Multiple Precision Arithmetic Library},
6.3.0 ed., 2023,
\url{https://gmplib.org/}.

\bibitem{Guy1983}
R.~K. Guy,
\emph{Don't try to solve these problems},
Amer.\ Math.\ Monthly \textbf{90} (1983), 35--41.

\bibitem{HuggableHacker2026}
Huggable Hacker,
\emph{Collatz Frontier: Distributed computational verification of the
  Collatz Conjecture beyond $2^{68}$},
GitHub repository, 2026,
\url{https://github.com/huggablehacker/Collatz-Frontier}.

\bibitem{Lagarias2003annotated}
J.~C. Lagarias,
\emph{The $3x{+}1$ problem: An annotated bibliography},
arXiv:math/0309224, 2003.

\bibitem{Lagarias1985}
J.~C. Lagarias,
\emph{The $3x{+}1$ problem and its generalizations},
Amer.\ Math.\ Monthly \textbf{92} (1985), no.~1, 3--23.

\bibitem{Matthews1984}
K.~R. Matthews and A.~M. Watts,
\emph{A generalization of Hasse's generalization of the Syracuse
  algorithm},
Acta Arith.\ \textbf{43} (1984), 167--175.

\bibitem{OliveiraeSilva2010}
T.~Oliveira e Silva,
\emph{Empirical verification of the $3x{+}1$ and related conjectures},
in \emph{The Ultimate Challenge: The $3x{+}1$ Problem}
(J.~C.~Lagarias, ed.),
Amer.\ Math.\ Soc., Providence, RI, 2010, pp.~189--207.

\bibitem{PythonDocs}
Python Software Foundation,
\emph{Python Language Reference, version 3.x},
\url{https://docs.python.org/3/}, 2023.

\bibitem{Roosendaal}
E.~Roosendaal,
\emph{$3x{+}1$ delay records},
\url{http://www.ericr.nl/wondrous/delrecs.html}, 2023
(accessed March 2026).

\bibitem{Shanks1993}
D.~Shanks,
\emph{Solved and Unsolved Problems in Number Theory},
4th ed., Chelsea Publishing, New York, 1993.

\bibitem{Tao2022}
T.~Tao,
\emph{Almost all Collatz orbits attain almost bounded values},
Forum Math.\ Pi \textbf{10} (2022), e12.
\doi{10.1017/fmp.2022.8}

\bibitem{Terras1976}
R.~Terras,
\emph{A stopping time problem on the positive integers},
Acta Arith.\ \textbf{30} (1976), 241--252.

\bibitem{Waitress}
P.~Cortesi,
\emph{Waitress: A production-quality pure-Python WSGI server},
\url{https://docs.pylonsproject.org/projects/waitress/}, 2023.

\end{thebibliography}

%=======================================================================
\appendix
\section{Complete Source Code}
\label{sec:code}
%=======================================================================

The following listings present the complete source code of the three
principal components.  All code is released under the MIT License.
The canonical, actively maintained version is available at
\url{https://github.com/huggablehacker/Collatz-Frontier}.

\subsection{\texttt{collatz\_coordinator.py}}

\begin{lstlisting}[caption={Coordinator: chunk dispatch, checkpointing,
                             and results collection.},
                   label={lst:coordinator}]
"""
Collatz Distributed Coordinator  -  Production Ready
Supports Gunicorn, Waitress, and Flask dev server (auto-detected).
"""
import json, logging, os, sys, threading, time
from pathlib import Path
from flask import Flask, request, jsonify

DEFAULT_START      = 2 ** 68
DEFAULT_CHUNK_SIZE = 500_000
DEFAULT_PORT       = 5555
RESULTS_FILE       = "collatz_distributed_fails.txt"
CHECKPOINT_FILE    = "collatz_coordinator_checkpoint.json"
STALE_TIMEOUT      = 600

log = logging.getLogger("werkzeug")
log.setLevel(logging.ERROR)
lock  = threading.Lock()
state = {
    "next_n": DEFAULT_START, "chunk_size": DEFAULT_CHUNK_SIZE,
    "chunks_issued": 0, "chunks_done": 0, "total_tested": 0,
    "fails": 0, "in_flight": {}, "started_at": time.time(),
}

def save_checkpoint():
    data = {k: v for k, v in state.items()
            if k not in ("in_flight", "started_at")}
    tmp = CHECKPOINT_FILE + ".tmp"
    with open(tmp, "w") as f: json.dump(data, f, indent=2,
                                         default=str)
    Path(tmp).replace(CHECKPOINT_FILE)

def load_checkpoint():
    p = Path(CHECKPOINT_FILE)
    if not p.exists(): return False
    try:
        with open(p) as f: saved = json.load(f)
        state.update(saved)
        state["in_flight"] = {}
        state["started_at"] = time.time()
        return True
    except Exception as e:
        print(f"  Warning: checkpoint unreadable ({e})")
        return False

def _startup(start_n=None, chunk_size=None):
    if start_n is None:
        start_n = int(os.environ.get("COLLATZ_START",
                                      DEFAULT_START))
    if chunk_size is None:
        chunk_size = int(os.environ.get("COLLATZ_CHUNK",
                                         DEFAULT_CHUNK_SIZE))
    resumed = load_checkpoint()
    if not resumed:
        n = start_n if start_n % 2 == 1 else start_n + 1
        state["next_n"] = n
        state["chunk_size"] = chunk_size
        with open(RESULTS_FILE, "w") as rf:
            rf.write("COLLATZ CONJECTURE -- DISTRIBUTED FAILS\n")
            rf.write(f"Started : {time.strftime('%Y-%m-%d %H:%M:%S')}\n")
            rf.write(f"Start n : 2^68 = {DEFAULT_START}\n")
            rf.write("-" * 100 + "\n")
    else:
        with open(RESULTS_FILE, "a") as rf:
            rf.write(f"\n--- Restarted "
                     f"{time.strftime('%Y-%m-%d %H:%M:%S')} "
                     f"(resuming n={state['next_n']:,}) ---\n")
    threading.Thread(target=_watchdog, daemon=True).start()
    return resumed

def _watchdog():
    while True:
        time.sleep(60); now = time.time()
        with lock:
            stale = [cid for cid, info in
                     list(state["in_flight"].items())
                     if now - info["issued_at"] > STALE_TIMEOUT]
            for cid in stale:
                info = state["in_flight"].pop(cid)
                if info["start"] < state["next_n"]:
                    state["next_n"] = info["start"]
                    state["chunks_issued"] -= 1

app = Flask(__name__)

@app.route("/chunk", methods=["GET"])
def get_chunk():
    worker_id = request.args.get("worker", "unknown")
    with lock:
        chunk_id = state["chunks_issued"]
        start = state["next_n"]
        if start % 2 == 0: start += 1
        end = start + state["chunk_size"] * 2 - 2
        state["next_n"] = end + 2
        state["chunks_issued"] += 1
        state["in_flight"][str(chunk_id)] = {
            "worker": worker_id, "start": start,
            "end": end, "issued_at": time.time()}
        save_checkpoint()
    return jsonify({"chunk_id": chunk_id, "start": start,
                    "end": end,
                    "chunk_size": state["chunk_size"]})

@app.route("/results", methods=["POST"])
def post_results():
    data      = request.get_json(force=True)
    chunk_id  = str(data.get("chunk_id"))
    worker_id = data.get("worker", "unknown")
    fails     = data.get("fails", [])
    count     = int(data.get("count", 0))
    elapsed   = float(data.get("elapsed_sec", 1))
    if fails:
        with open(RESULTS_FILE, "a", buffering=1) as rf:
            for r in fails:
                ts = time.strftime("%Y-%m-%d %H:%M:%S")
                rf.write(f"n={r['n']:<30} | FAIL | "
                         f"steps={r['steps']:<10} | "
                         f"peak={r['peak']} | "
                         f"found={ts} | worker={worker_id}\n")
    with lock:
        state["chunks_done"]  += 1
        state["total_tested"] += count
        state["fails"]        += len(fails)
        state["in_flight"].pop(chunk_id, None)
        save_checkpoint()
    return jsonify({"status": "ok"})

@app.route("/status", methods=["GET"])
def status():
    with lock:
        elapsed = time.time() - state["started_at"]
        rate = (state["total_tested"] / elapsed
                if elapsed > 0 else 0)
        workers = list({v["worker"]
                        for v in state["in_flight"].values()})
    lines = ["=" * 60, " COLLATZ COORDINATOR", "=" * 60,
             f" Next n : {state['next_n']}",
             f" Tested : {state['total_tested']:,}",
             f" Rate   : {rate:,.0f} odd/sec",
             f" Workers: {workers}",
             f" FAILs  : {state['fails']}", "=" * 60]
    return "<pre>" + "\n".join(lines) + "</pre>"

def post_fork(server, worker): _startup()

_gunicorn = ("gunicorn" in
             os.environ.get("SERVER_SOFTWARE", "") or
             any("gunicorn" in a for a in sys.argv))
if _gunicorn: _startup()

if __name__ == "__main__":
    import argparse
    p = argparse.ArgumentParser()
    p.add_argument("--start", type=int, default=DEFAULT_START)
    p.add_argument("--chunk", type=int,
                   default=DEFAULT_CHUNK_SIZE)
    p.add_argument("--port",  type=int, default=DEFAULT_PORT)
    args = p.parse_args()
    _startup(start_n=args.start, chunk_size=args.chunk)
    try:
        from waitress import serve
        print("  Server: Waitress (production)")
        serve(app, host="0.0.0.0", port=args.port, threads=8)
    except ImportError:
        print("  Server: Flask dev server")
        app.run(host="0.0.0.0", port=args.port, threaded=True)
\end{lstlisting}

\subsection{\texttt{collatz\_worker.py}}

\begin{lstlisting}[caption={Worker: multiprocessing dispatch and
                             Syracuse kernel.},
                   label={lst:worker}]
"""
Collatz Distributed Worker  -  Optimized
Multiprocessing, Syracuse steps, odd-only, early-exit.
"""
import argparse, socket, sys, time
import requests, multiprocessing
from multiprocessing import Pool

MAX_STEPS = 10_000_000
THRESHOLD = 2 ** 68

def process_subchunk(args):
    start, end, threshold, max_steps = args
    try:
        import gmpy2; gmpy2.get_context().precision = 256
        def mpz(n): return gmpy2.mpz(n)
    except ImportError:
        def mpz(n): return n
    thr = mpz(threshold); fails = []; count = 0; n = start
    while n <= end:
        steps = 0; max_val = n; cur = mpz(n)
        while cur >= thr:
            if steps > max_steps:
                fails.append({"n": n, "steps": steps,
                               "peak": int(max_val)})
                break
            x = 3 * cur + 1
            k = (x & -x).bit_length() - 1
            cur = x >> k; steps += k + 1
            if cur > max_val: max_val = cur
        count += 1; n += 2
    return fails, count

def run_worker(coordinator_url, worker_id, num_cores):
    session = requests.Session()
    session.headers["Content-Type"] = "application/json"
    with Pool(processes=num_cores) as pool:
        while True:
            try:
                resp = session.get(
                    f"{coordinator_url}/chunk",
                    params={"worker": worker_id}, timeout=30)
                resp.raise_for_status()
                chunk = resp.json()
            except Exception as e:
                print(f"  [{worker_id}] Chunk error: {e}")
                time.sleep(5); continue
            chunk_id = chunk["chunk_id"]
            start = int(chunk["start"]); end = int(chunk["end"])
            total = (end - start) // 2 + 1
            per   = max(1, total // num_cores)
            subs  = []; s = start
            for i in range(num_cores):
                e2 = s + (per - 1) * 2
                if i == num_cores - 1: e2 = end
                subs.append((s, e2, THRESHOLD, MAX_STEPS))
                s = e2 + 2
                if s > end: break
            t0 = time.time()
            results = pool.map(process_subchunk, subs)
            elapsed = time.time() - t0
            all_fails = []; count = 0
            for f, c in results: all_fails.extend(f); count += c
            try:
                session.post(
                    f"{coordinator_url}/results",
                    json={"chunk_id": chunk_id,
                          "worker": worker_id,
                          "fails": all_fails, "count": count,
                          "elapsed_sec": elapsed},
                    timeout=60).raise_for_status()
            except Exception as e:
                print(f"  [{worker_id}] Post error: {e}")
            rate = count / elapsed if elapsed > 0 else 0
            print(f"  [{worker_id}] chunk={chunk_id} "
                  f"| {count:,} | {rate:,.0f}/sec "
                  f"| fails={len(all_fails)}")

if __name__ == "__main__":
    multiprocessing.freeze_support()
    p = argparse.ArgumentParser()
    p.add_argument("--coordinator", required=True)
    p.add_argument("--name", default=socket.gethostname())
    p.add_argument("--cores", type=int,
                   default=multiprocessing.cpu_count())
    args = p.parse_args()
    try:
        run_worker(args.coordinator.rstrip("/"),
                   args.name, args.cores)
    except KeyboardInterrupt:
        print(f"\n  [{args.name}] Stopped.")
\end{lstlisting}

\subsection{\texttt{collatz\_merge.py}}

\begin{lstlisting}[caption={Results merger: combines distributed FAIL
                             files into a single sorted report.},
                   label={lst:merge}]
"""Collatz Results Merger."""
import argparse, re, time
from pathlib import Path

FAIL_RE = re.compile(
    r"n=\s*(\d+)\s*\|.*?FAIL.*?\|\s*steps=\s*(\d+)"
    r"\s*\|\s*peak=(\d+)")

def parse_file(path):
    records = []
    try:
        with open(path) as f:
            for line in f:
                m = FAIL_RE.search(line)
                if m:
                    records.append({
                        "n":     int(m.group(1)),
                        "steps": int(m.group(2)),
                        "peak":  int(m.group(3)),
                        "raw":   line.rstrip()})
    except FileNotFoundError: pass
    return records

if __name__ == "__main__":
    p = argparse.ArgumentParser()
    p.add_argument("--main",
                   default="collatz_distributed_fails.txt")
    p.add_argument("--out",
                   default="collatz_merged_fails.txt")
    args = p.parse_args()
    files = ([Path(args.main)] +
             sorted(Path(".").glob("collatz_local_fails_*.txt")))
    records = []
    for f in files: records.extend(parse_file(f))
    records.sort(key=lambda r: r["n"])
    records = list({r["n"]: r for r in records}.values())
    with open(args.out, "w") as f:
        ts = time.strftime('%Y-%m-%d %H:%M:%S')
        f.write(f"COLLATZ MERGED FAILS -- Generated: {ts}\n")
        f.write(f"Total FAILs: {len(records)}\n")
        f.write("-" * 100 + "\n")
        for r in records: f.write(r["raw"] + "\n")
    print(f"  {len(records)} FAILs written to {args.out}")
\end{lstlisting}

\end{document}
